\documentclass{article}
\usepackage{/home/user1/code/tex/sty}
\usepackage{amsfonts}
\usepackage{amsmath}
\usepackage{geometry}
\usepackage{graphicx}
\usepackage{hyperref}
\usepackage{enumitem}
\setlength{\parindent}{0pt}
\geometry{top=1in,bottom=1in,left=1in,right=1in}
\begin{document}
\begin{center}
\textbf{\Large{Woche 5 Hausaufgaben}}\\~\\
\begin{tabular}{l l}
Student&Agustin Romero\\
Email&ar1@stanford.edu\\
Instructor&Patric Di Dio De Marco\\
Course&GERLANG 3\\
Institution&Stanford University\\
Date&\today\\
\end{tabular}
\end{center}
\section*{Hausafgaben Woche 5 Tag 1}
% AB 196-3, 196-4, 197-5, 197-6
\subsection*{196-3}
\listalph{\item Werkzeugmacher sind Handwerker, die Werkzeug herstellen.
\item Sicherheitsvorschriften sind Regeln für die Sicherheit, die auf jeder Baustelle beachtet werden müssen.
\item Ein Steinbruch ist ein Ort, an dem Stine für Bauprojekte gebrochen werden.
\item Im Mittelalter gab es auf den Baustellen oft Unfälle, bei denen viele Menschen verletzt wurden.
\item Auf der Baustelle in Meßkirch dürfen nur Seile verwendet werden, die aus Stahl machten.
\item Zimmerleute sind Handwerker, die aus Holz Dächer und andere Bauteile bauen.
}
\subsection*{196-4}
\listalph{
\item Die Leitern wurden weggeräumt.
\item Die Bänke wurden aufgestellt.
\item Das Fenster wurde eingesetzt.
\item Der Boden wurden geputzt.
\item Die Bilder wurden aufgehängt.
\item Die Tiere wurden aus der Kirche gebracht.
\item Die Werkzeuge wurden weggeräumt.
\item Die Steine wurden aus der Kirche getragen.
\item Die Wände wurden geschmückt.
}
\subsection*{197-5}
\listnum{\item e. Die Schuldirektorin Manuela Lenz sagte das Schulfest ab.
\item b. Ein Hochwasser zerstörte die Brüche über den Leiterbach.
\item f. Dr. Müller vom Stadtamt verbot das Fahrradfahren im Stadtpark.
\item a. Ein Sturm beschädigte das Dach der Kirche.
\item d. Der Besitzer Walter Kern verkaufte das Grundstück am See.
\item c. Ein Blitz starb tötete fünf Kühe.
}
\subsection*{197-6}
\listnum{\item Das Schulfest wurde von der Schuldirektorin Manuela Lenz abgesagt.
\item Die Brüche über den Leiterbach wurde von dem Hochwasser wurde zerstört.
\item Das Fahrradfahren im Stadpark wurde von Dr. Müller verboten.
\item Der Kirche wurde von dem Sturm beschädigt.
\item Das Grundstück am See wurde von der Besitzer Walter Kern verkauft.
\item Funf Kühe wurde von ein Blitz gestorben.
}
\section*{Hausafgaben Woche 5 Tag 2}
% AB 197-7, 198-5, 199-6
\subsection*{197-7}
\listalph{
\item Ein Bauarbeiter hatte bei einem Unfall schwer verletzt.
\item Das neues Autobahnteilstück hatte geöffnet.
\item Zwei Jungendliche Kühe wurde bei Autounfall hatten getotet.
\item Das Eisenbahnbrucke nach sechs Monaten Bauzeit hat fertiggestellt.
\item Vier kleine Wasserkraftwerke in Betrieb hatten genommen.
\item Das Jungendzentrum hat geschlossen.
}
\subsection*{198-5}
\listalph{
\item Manuela hatte ein Problem mit ihrem Scanner. Falsch.
\item Bei ihrem ersten Anruf wurde sie zweimal verbunden. Richtig.
\item Sie ärgert sich, weil sie wie eine Anfängerin behandelt wurde. Richtig. 
\item Manuela wurde von einer Hotline-Mitarbeiterin zurückgerufen. Richtig.
\item Manuela musste noch einmal ihr Problem beschreiben. Richtig.
\item Das Problem konnte von der Hotline-Mitarbeiterin gelöst werden. Falsch.
}
\subsection*{199-6}
\listnum{
\item Ich bin zweimal bei meinem ersten Anruf verbunden worden.
\item Ich bin wie eine Anfängerin mich geärgert worden.
\item Ich bin von einer Hotline-Mitarbeiterin zurückgeruft worden. 
\item Mein Problem ist nicht von der Hotline-Mitarbeiterin gelöst worden.
}
\section*{Hausafgaben Woche 5 Tag 3}
% AB 197-7, 198-5, 199-6
Keine Hausaufgaben von dem Arbeitsbuch wurde gebietet.
\section*{Hausafgaben Woche 5 Tag 4}
% AB S.200 (3, 5), S.201 (6), S.202 Schreibwerkstatt (3)
\subsection*{200-3}
\listpunkt{\item Ist der Aufnahmtest diesen Donnerstag? \listpunkt{\item Nein, er ist voraussichtlich nächsten Montag.}
\item Letztes Jahr haben wir in Italien Urlaub gemacht. \listpunkt{\item Und wohin fahrt ihr nächstes Jahr?}
\item Spielen wir diese Woche Tennis? \listpunkt{\item Da kann ich nicht, aber nächste Woche kann ich.}
\item Dieser Winter ist ziemlich warm. \listpunkt{\item Stimmt. Ich hoffe, wir Skifahrer haben nächsten Winter mehr Glück.}
\item Wo fahrt ihr letztes Wochenende? \listpunkt{\item Nirgendwo, wir haben uns zu Hause ausgeruht.}}
\subsection*{200-5}
\listpunkt{\item Bis wann willst du denn noch an deinen Gedichten arbeiten? \listpunkt{\item Bis sie perfekt sind. Sprachliche Intelligenz.}
\item Seit wann kann denn sein Sohn schwimmen? \listpunkt{\item Seit wir letzten Sommer am Bodensee Urlaub gemacht haben. Körperliche Intelligenz.}
\item Bis wann müssen wir die Pläne fertig zeichnen? \listpunkt{\item Bis Freitag, dann werden sie nach Münschen geschickt. Räumliche Intelligenz.}
\item Seit wann magst du Sudoku? \listpunkt{\item Seit mir Caro kürzlich gezeigt hat, wie das funktionert. Mathematische Intelligenz.}
\item Seit wann spielst du Klavier? \listpunkt{\item Schon seit meiner Kindheit. Ich hatte shon damals Klavierunterricht. Musikalische Intelligenz.}
\item Seit wann willst du denn noch in der Sonne liegen und nichts tun? \listpunkt{\item Bis ich mich so richtig enspannt habe. Personale Intelligenz.}}
\subsection*{201-6}
\listalph{\item Seit Maria ihre Freunde in Shanghai getroffen hat, sind zwei Monate vergangen.
\item Seit Maria nach Asien geflogen hat, sind drei Monate vergangen.
\item Seit Maria aus Asien zurückgekommen hat, sind zwei Woche vergangen.
\item Bis Maria für die ersten Prüfungen lernen müssen, sind es noch drei Monate.
\item Bis Maria in die WG einziehen, sind es noch drei Wochen.
\item Bis das Studium von Maria beginnen werden, sind es noch zwei Monate.}
\subsection*{202-3}
Hier eine kleine Zeitungsmeldung aus Sonnental wird beschreibt.

Zuerst hat ein Restaurant der Inseketenessen wurde geöffnet. Wahrscheinlich der Einwohner sind für Insektessen, und meinen dass es gut geschmekt, denn an dem ersten Eröffnungsfeier war das Restuarant voll. Andererseitz könnten Menschen widersprechen, zum Beispiel sie würden gerne Obst und Fleisch essen, trotz Insekten sind billiger als Fleisch.\\~\\
Zunächst wird der neue Gesetz über der Verkehr des Stadtzentrums angefangen. Der neue Erklarung mitteilt, dass nur Elektrofahrzeugen durfen buchen auf dem Stadtzentrum, also kein Ottomotor und kein Dieselmotor werden erlauben seit nächsten Monat. Mansche Menschen können gegensprechen über dieses neues Gesetz, insbesondere der Einwohner auf dem Stadtzentrums, die braucht schon ein Auto. Andererseits haben Menschen wahrscheinlich das Problem, dass die Autos in dem Stadt sind unregiert, zum Beispiel manche Autos braucht keine Katalysator, also die Luftqualität wird schelchter werden, falls dies Gesetz wird nicht aufpassen.\\~\\
Zunächst wird neues Robot entwickelt, um alte Menschen von Pflegerheimen zu sorgen. Alle Staat haben die bestimmte Problem, um der alter Generation des Menschen zu pflegen. Denn nicht alle alter Menschen können auf unseren Heimat leben, müssen manche auf Pflegerheim wohnen. Tatsächlich ist, dass viele Aktivität auf Pflegerheimen können automatisch erfolgen, zum Beispiel der reinen des Flurs, der erzeugen des Nahrungsmittels und so weiter. Aber die Mehrheit des Pflegewerk sollen von studierenden Pfleger, denn viele Aktivitäten sind mehr kompliziert als was Roboten können brauchen, zum Beispiel die Pillen auswählen und Gymnastik zu leiten. Es ist aber möglich, dass mehr Aktivität auf alle Heimaten werden brauchen von automatisch Geräten.\\~\\
Endlich haben uns die 
\section*{Hausafgaben Woche 5 Tag 5}
Keine Hausaufgaben von dem Arbeitsbuch wurde gebietet.
\end{document}
